\chapter*{Preface: Bring attention back to the research}
\addcontentsline{toc}{chapter}{Preface}
\markboth{Preface}{Preface}

We often leave figures until the final stages of writing a paper. The data have been analyzed, the methods are documented, and the conclusions seem clear. All that remains is to make the figures look good. Once we begin, problems appear: one image cannot hold everything, two arrows invite different interpretations, a summary hides individual variation, and text that looked clear on a monitor becomes unreadable at its final size.

These difficulties are closely connected to the research. Making a figure requires us to decide again what matters most, which elements are evidence, which are interpretation, and in what order the reader should encounter them. Sometimes a figure is unclear because its layout has not been resolved. Sometimes the question it should answer has not yet been stated precisely.

AI lowers the effort needed to begin. We can describe a scene in words, suggest a composition with a sketch, explain a revision using an existing image, or give data to a visualization tool and explore ways to present them. Tasks that once required many separate operations can now start with a short instruction. That creates more room to compare alternatives and refine the explanation.

A polished image can also conceal mistakes. An incorrect relationship can look convincing. A tidy layout can distract from a missing experimental condition. An unsupported conclusion can acquire a professional-looking visual form. We therefore need to move from being able to generate a figure to being able to explain, inspect, and deliver it.

This book follows that process. You will begin with a drawing brief that states the purpose of the figure and the question it should answer. You will then choose a workflow for a data plot, a methods diagram, a mechanism illustration, an experimental image, or a graphical abstract. We will discuss composition, color, labels, and legends, as well as revision, export, and preservation of source materials. Each part explains what AI can assist with and what must be decided from the research record.

SciDraw AI is the practical platform used throughout the book. A common platform lets readers continue working on the same example without repeatedly learning a new interface. Instructions are based on actual interfaces and available operations. Where a workflow requires a download, an upload, or a manual addition, that step is stated directly. Features enter a tutorial after they have been implemented and checked.

You do not need to master illustration before starting. More useful preparation is a study you understand, a complete set of source materials, and a willingness to inspect each important element. Experienced readers can enter at a specific problem: retaining the study design in a data plot, reducing ambiguity in a mechanism diagram, or arranging several panels into a readable argument.

Exercises distinguish simulated materials from real research. Simulated data allow us to share a reproducible task and discuss mistakes without presenting a teaching result as a scientific discovery. Real cases require reliable sources and any necessary permission. The aim is a method that transfers to your own work: decide what to communicate, use the tools, and preserve important decisions together with their basis.

In the next chapter, focus on one concrete question: what should this figure help the reader understand?
